\documentclass[11pt,a4paper]{article}
\usepackage{fontspec,amsmath,amssymb,amsthm,microtype}
\setmainfont{lmroman10-regular.otf}[BoldFont=lmroman10-bold.otf,ItalicFont=lmroman10-italic.otf,BoldItalicFont=lmroman10-bolditalic.otf]
\setmonofont{lmmono10-regular.otf}[Scale=MatchLowercase]
\usepackage[margin=25mm]{geometry}
\usepackage[unicode,hidelinks]{hyperref}
\setlength{\emergencystretch}{3em}
\setlength{\parskip}{0.4em}
\setlength{\parindent}{0pt}
\allowdisplaybreaks[2]
\raggedbottom
\title{\texorpdfstring{Prove the actual Fisher dissipation and sharp frozen-\allowbreak{}gradient forcing}{Prove the actual Fisher dissipation and sharp frozen-gradient forcing}}
\author{LangevinLean mathematical exposition}
\date{}
\begin{document}\maketitle
\noindent Source module: \texttt{Unadjusted\allowbreak{}Langevin/\allowbreak{}LSI/\allowbreak{}Gaussian\allowbreak{}Entropy\allowbreak{}Production.\allowbreak{}lean}. The accompanying HTML pairs this mathematical exposition with the exact Lean statements and records the scope of the formalization.


The actual Gaussian entropy chain rule gives an almost-everywhere derivative represented by a flux on the independent initial/\allowbreak{}noise probability space. Averaging the genuine Gaussian component score identifies this flux with \(-I_t+\mathbb E\langle\nabla U(Y_t)-\nabla U(X),\nabla\log r_t(Y_t)\rangle\). Young's inequality absorbs a quarter of \(I_t\); the already proved exact Gaussian displacement estimate bounds the remaining error. The result is the unconditional analytic production inequality \(H'_t\le-\tfrac34I_t+t^2L^2G_0+2tdL^2\) almost everywhere at positive times, with no interpolation regularity or PDE data hypotheses.


\section{\texorpdfstring{Justify inner-\allowbreak{}product integrability from two L2 fields}{Justify inner-product integrability from two L2 fields}}


For vector fields \(F,G\in L^2(P)\), Hölder's inequality puts \(\|F\|\|G\|\) in \(L^1(P)\). Since \(|\langle F,G\rangle|\le\|F\|\|G\|\), their inner product is integrable.



This private Lean helper applies \texttt{Mem\allowbreak{}Lp.\allowbreak{}mul} to the two norm fields with resulting exponent one, converts it to integrability, and uses \texttt{Integrable.\allowbreak{}mono'} with \texttt{norm\_\allowbreak{}inner\_\allowbreak{}le\_\allowbreak{}norm}. Measurability follows from the two input fields. The helper is used for the initial gradient, current gradient, Gaussian score, and relative-score pairings before splitting or recombining their integrals.



{\raggedright\noindent Lean declarations: \texttt{production\_\allowbreak{}integrable\_\allowbreak{}inner\_\allowbreak{}of\_\allowbreak{}mem\allowbreak{}Lp\_\allowbreak{}two}.\par}

\section{\texorpdfstring{Identify relative-\allowbreak{}score energy on the independent noise space with actual Fisher information}{Identify relative-score energy on the independent noise space with actual Fisher information}}


Let \(S_t(X,Z)=\nabla\log r_t(Y_t)\) and \(P=\mu\otimes\mathcal N(0,I)\). The exact update-law identity gives

\[\mathbb E_P\|S_t\|^2=\int\|\nabla\log r_t\|^2\,d\rho_t=I_t.\]

Lean first obtains integrability of the actual state-space relative score squared. \texttt{integral\_\allowbreak{}map} transports this integral through the measurable Gaussian update. The pointwise equality \texttt{hpoint} identifies the jointly measurable score field used along the update with the actual gradient of the log relative density. Finally the definition of \texttt{relative\allowbreak{}Fisher}, the actual tilted-law equality, and the ordinary/\allowbreak{}nonnegative integral conversion identify the finite real Fisher value. No energy equality is supplied as a hypothesis.



{\raggedright\noindent Lean declarations: \texttt{interpolation\allowbreak{}Score\allowbreak{}Along\allowbreak{}Update\_\allowbreak{}energy\_\allowbreak{}eq\_\allowbreak{}fisher}.\par}

\section{\texorpdfstring{Match the measurable state-\allowbreak{}space score field to the same finite Fisher value}{Match the measurable state-space score field to the same finite Fisher value}}


The jointly measurable field used in time-space estimates agrees with the actual relative score at each positive time. Consequently

\[\int\|\mathrm{ScoreField}_t(y)\|^2\,d\rho_t(y)=I_t.\]

Lean converts the previously proved actual-score \(L^2\) result to \texttt{Mem\allowbreak{}Lp} for this explicit field using \texttt{interpolation\allowbreak{}Relative\allowbreak{}Score\allowbreak{}Field\_\allowbreak{}eq}. It then maps the squared norm to the independent update space and invokes the preceding Fisher-energy equality. This bridge allows the LSI inequality, the time-space domination bounds, and the final production estimate to use interchangeable but formally identified score representations.



{\raggedright\noindent Lean declarations: \texttt{interpolation\allowbreak{}Relative\allowbreak{}Score\allowbreak{}Field\_\allowbreak{}energy\_\allowbreak{}eq\_\allowbreak{}fisher}.\par}

\section{\texorpdfstring{Integrate Young absorption against the actual frozen-\allowbreak{}gradient discrepancy}{Integrate Young absorption against the actual frozen-gradient discrepancy}}


Write \(D_t=\nabla U(Y_t)-\nabla U(X)\). The pointwise inequality

\[\langle D_t,S_t\rangle\le\frac14\|S_t\|^2+\|D_t\|^2\]

is integrable because both squared norms have already been proved integrable. Their expectations satisfy \(\mathbb E\|S_t\|^2=I_t\) and

\[\mathbb E\|D_t\|^2\le L^2\bigl(t^2G_0+2td\bigr),\qquad G_0=\mathbb E\|\nabla U(X)\|^2.\]

Lean constructs discrepancy measurability as \texttt{h\allowbreak{}Dm}, score measurability as \texttt{h\allowbreak{}Sm}, and the two energy integrabilities as \texttt{h\allowbreak{}Ei} and \texttt{h\allowbreak{}Si}. The weighted-inner-product helper with weight one proves cross-term integrability before applying the integrated Young inequality. \texttt{h\allowbreak{}F} replaces joint score energy by finite actual Fisher information; \texttt{h\allowbreak{}Eb} is the previously proved sharp Gaussian discrepancy bound. Combining them gives exactly \(\mathbb E\langle D_t,S_t\rangle\le I_t/4+t^2L^2G_0+2tdL^2\).



{\raggedright\noindent Lean declarations: \texttt{interpolation\_\allowbreak{}error\_\allowbreak{}score\_\allowbreak{}young}.\par}

\section{\texorpdfstring{Convert the actual Gaussian entropy flux into Fisher dissipation and gradient error}{Convert the actual Gaussian entropy flux into Fisher dissipation and gradient error}}


The raw entropy flux is \(-\mathbb E\langle g(X)-\sqrt{2t}Z/(2t),S_t\rangle\). The averaged Gaussian score identity gives

\[\mathbb E\left\langle-\frac{\sqrt{2t}}{2t}Z,S_t\right\rangle=\int\langle\nabla\log p_t,\nabla\log r_t\rangle\,d\rho_t.\]

Since \(\nabla\log r_t=\nabla U+\nabla\log p_t\), the right side equals \(I_t-\mathbb E\langle\nabla U(Y_t),S_t\rangle\). Substituting and collecting the two gradient terms yields

\[\mathrm{RawFlux}_t=-I_t+\mathbb E\langle\nabla U(Y_t)-\nabla U(X),S_t\rangle.\]

The Lean proof names the pulled-back initial gradient \texttt{GX}, current gradient \texttt{GY}, Gaussian component score \texttt{N}, and relative score \texttt{S}. All four are in \(L^2(P)\), so \texttt{h\allowbreak{}X}, \texttt{h\allowbreak{}G}, \texttt{h\allowbreak{}Noise}, and \texttt{h\allowbreak{}Out} establish the integrability required for every sum and difference. The map identities \texttt{h\allowbreak{}Map\allowbreak{}G} and \texttt{h\allowbreak{}Map\allowbreak{}S} identify output-space integrals with the actual independent-noise integrals. \texttt{h\allowbreak{}Leb} is the pointwise score decomposition, \texttt{h\allowbreak{}Avg'} is its integrated Gaussian-score consequence, and \texttt{h\allowbreak{}Raw}/\allowbreak{}\texttt{h\allowbreak{}Err} perform the two final algebraic expansions. This connects the analytic entropy flux to actual Fisher information without assuming a conditional drift equation.



{\raggedright\noindent Lean declarations: \texttt{interpolation\allowbreak{}Raw\allowbreak{}Entropy\allowbreak{}Flux\_\allowbreak{}eq\_\allowbreak{}fisher\_\allowbreak{}error}.\par}

\section{\texorpdfstring{Retain three quarters of actual Fisher dissipation}{Retain three quarters of actual Fisher dissipation}}


Insert the integrated Young estimate into the exact entropy-flux decomposition:

\[\mathrm{RawFlux}_t=-I_t+\mathbb E\langle D_t,S_t\rangle\le-\frac34I_t+t^2L^2G_0+2tdL^2.\]

The Lean proof rewrites the flux with \texttt{interpolation\allowbreak{}Raw\allowbreak{}Entropy\allowbreak{}Flux\_\allowbreak{}eq\_\allowbreak{}fisher\_\allowbreak{}error}, obtains the already justified numerical cross-term bound, and uses linear arithmetic. The target's probability normalization and finite initial gradient energy suffice for this pointwise flux estimate. LSI is needed later to relate Fisher information to entropy, rather than to justify Young absorption.



{\raggedright\noindent Lean declarations: \texttt{interpolation\allowbreak{}Raw\allowbreak{}Entropy\allowbreak{}Flux\_\allowbreak{}le}.\par}

\section{\texorpdfstring{Obtain the genuine positive-\allowbreak{}time entropy derivative inequality on one common full-\allowbreak{}measure set}{Obtain the genuine positive-time entropy derivative inequality on one common full-measure set}}


The actual KL entropy is absolutely continuous on every compact interval inside \((0,\infty)\), and its derivative equals the actual raw flux almost everywhere there. To use this fact on an entire finite range \((0,T)\), choose the countable intervals

\[\left(\frac1{n+1},T+1\right),\qquad n=0,1,2,\ldots.\]

Intersect their full-measure derivative sets. Every \(t\in(0,T)\) belongs to at least one such interval because \(1/(n+1)\to0\). At that time the actual derivative equals the raw flux, so the pointwise flux bound gives

\[H'_t\le-\frac34 I_t+t^2L^2G_0+2tdL^2\quad\text{for almost every }t\in(0,T).\]

Lean's \texttt{hderiv} applies the fully proved local entropy chain rule for each interval; \texttt{ae\_\allowbreak{}all\_\allowbreak{}iff} makes their countable intersection. \texttt{exists\_\allowbreak{}nat\_\allowbreak{}one\_\allowbreak{}div\_\allowbreak{}lt} chooses an interval containing the current positive time. \texttt{hd.\allowbreak{}deriv} identifies the ordinary derivative, and \texttt{interpolation\allowbreak{}Raw\allowbreak{}Entropy\allowbreak{}Flux\_\allowbreak{}le} supplies the numerical bound. This argument avoids asserting a derivative at time zero, where the initial law need not have a smooth positive density.



{\raggedright\noindent Lean declarations: \texttt{interpolation\allowbreak{}Entropy\_\allowbreak{}derivative\_\allowbreak{}bound\_\allowbreak{}ae}.\par}

\end{document}
